% !TeX program = xelatex
% 本文件只填写简历内容；颜色、字号、间距等排版设置均在 simplecnresume.cls 中。
% 下方姓名、联系方式、单位、项目、论文和奖项均为虚构示例，请直接替换。
\documentclass{simplecnresume}

% 推荐投递风格：经典深蓝、纯色章节线、紧凑竖线分隔符。
% 重点突出信息层级和正文可读性；其他预设及自定义方法见 README。
\resumesetup{
  color=ResumeClassic,
  section-line=solid,
  entry-separator=bar
}

\begin{document}

% 用法：\resumeheader[照片路径]{姓名}{联系信息}
% 不填写可选的照片路径时，将自动显示通用照片占位框。
\resumeheader[images/placeholder-id-photo-transparent.png]{示例姓名}{%
  \resumecontact{\faPhone}{联系电话：138-0000-0000}
  \resumecontact{\faEnvelope}{邮箱：\href{mailto:hello@example.com}{hello@example.com}}\\
  \resumecontact{\faGlobe}{技术博客：\href{https://example.com}{https://example.com}}
}

\resumesection{个人简介}
\begin{resumeitems}
  \item \textbf{个人概况：}计算机相关专业毕业生，关注大语言模型应用、数据分析与机器学习工程，具备从需求分析、方案设计到服务部署的完整项目实践。
  \item \textbf{能力亮点：}熟悉 Python、PyTorch、RAG、Agent 与常用工程化工具，能够将实验原型整理为可维护的业务服务。
\end{resumeitems}

\resumesection{工作经历}
% 投递版优先使用纯文字公司栏，减少装饰干扰并保持黑白打印清晰。
% 如需作品展示风格，也可换用带 Logo 的 \resumebanner，示例见 examples。
\resumeentry[black][bar]{星河智能科技有限公司}{AI 应用工程师}{2024.07 -- 至今}
\begin{resumeitems}
  \item 负责企业知识检索、智能问答与数据分析类应用的算法设计、效果评估及服务化部署。
\end{resumeitems}

\resumeentry[black][bar]{云杉数据实验室}{算法工程师（实习）}{2023.07 -- 2024.03}
\begin{resumeitems}
  \item 参与文本分类与信息抽取模型开发，完成数据清洗、实验设计、误差分析和接口封装。
\end{resumeitems}

\resumesection{项目经历}
% 项目标题使用主题色；第一个可选参数支持 black、ResumePrimary 或自定义颜色。
\resumeentry[ResumePrimary][bar]{企业知识库问答平台}{RAG / 大语言模型应用}{2024.08 -- 2025.01}
\begin{resumeitems}
  \item \textbf{项目背景：}面向内部制度、产品手册和常见问题构建可追溯的知识问答平台，减少人工检索成本。
  \item \textbf{核心技术：}Python、FastAPI、向量检索、关键词检索、Rerank、Docker。
  \item \textbf{主要工作：}
    \begin{resumedetails}
      \item \textbf{数据管线：}\resumecircled{1}设计文档解析、分块、索引和增量更新流程；\resumecircled{2}统一多格式文档的元数据结构。
      \item \textbf{检索优化：}\resumeparen{1}实现混合检索与重排序链路；\resumeparen{2}通过查询改写和引用校验改善复杂问题的回答质量。
      \item \textbf{效果评测：}\resumerightparen{1}建立事实性、完整性和引用正确率评测集；\resumerightparen{2}将回答准确率由 \resumehighlight{78\% 提升至 91\%}。
    \end{resumedetails}
\end{resumeitems}

\resumeprojectdivider
\resumeentry[ResumePrimary][bar]{多智能体数据采集助手}{Agent / 浏览器自动化}{2024.03 -- 2024.07}
\begin{resumeitems}
  \item \textbf{项目背景：}将自然语言采集需求转换为可执行步骤，辅助非技术用户完成公开网页数据整理。
  \item \textbf{核心技术：}LangGraph、ReAct、Playwright、Function Calling、结构化输出。
  \item \textbf{主要工作：}
    \begin{resumedetails}
      \item \textbf{Agent 架构：}设计规划、执行、页面分析和结果校验等角色，使用共享状态管理任务进度与异常信息。
      \item \textbf{稳定性设计：}增加限次重试、页面状态检测和断点恢复机制，使示例任务成功率达到 \resumehighlight{94\%}。
    \end{resumedetails}
\end{resumeitems}

\resumesection{个人技能}
\begin{resumeitems}
  \item \textbf{编程与框架：}Python、SQL、PyTorch、FastAPI、Git、Docker。
  \item \textbf{大模型应用：}Prompt Engineering、RAG、Agent、Function Calling、Text-to-SQL、模型评测。
  \item \textbf{机器学习：}分类与回归、时间序列分析、Transformer、实验设计与误差分析。
  \item \textbf{数据处理：}Pandas、NumPy、MySQL、数据清洗、特征工程与可视化。
\end{resumeitems}

\resumesection{教育背景}
\resumeentry[black]{海棠理工大学}{计算机科学与技术 \quad 本科}{2020.09 -- 2024.06}
\begin{resumeitems}
  \item \textbf{GPA：}3.7/4.0（专业前 10\%）。
  \item \textbf{荣誉奖项：}校级优秀学生奖学金、优秀毕业生、优秀毕业设计。
  \item \textbf{主修课程：}数据结构、机器学习、数据库系统、自然语言处理。
\end{resumeitems}

\resumesection{学科竞赛}
\begin{resumeawards}
  \resumeaward{星桥杯数据科学挑战赛全国一等奖}
  \resumeaward{全国高校算法应用示例赛二等奖}
  \resumeaward{青年人工智能创新竞赛银奖}
  \resumeaward{校级数学建模挑战赛一等奖}
\end{resumeawards}

\end{document}
